% ថ្នាក់ទី ១១ (វិទ្យាសាស្ត្រ) · ជំពូកទី ១ · មេរៀនទី ២ · សៀវភៅទំព័រ ៨–១៣
\chapter{ស៊ីតនព្វន្ត}

\begin{objBox}[វត្ថុបំណង]
\begin{enumerate}[1]
  \item កំណត់និយមន័យនិងផលសងរួមនៃស៊ីតនព្វន្ត
  \item កំណត់តួទី $n$ នៃស៊ីតនព្វន្ត
  \item ផលបូកតួស្មើចម្ងាយពីតួចុង
  \item គណនាផលបូកតួនៃស៊ីតនព្វន្ត ។
\end{enumerate}
\end{objBox}

\section{និយមន័យស៊ីតនព្វន្ត}

\begin{example}[ឧទាហរណ៍]
គេមានស៊ីត $1, 2, 3, 4, 5, \ldots$ ។ គេបាន
\[
\begin{aligned}
a_1 &= 1,\ a_2 = 2 = 1 + 1 = a_1 + 1,\ a_3 = 3 = 2 + 1\\
&= a_2 + 1,
\end{aligned}
\]
\[
\begin{gathered}
\begin{aligned}
a_4 &= 4 = 3 + 1 = a_3 + 1,\ a_5 = 5 = 4 + 1\\
&= a_4 + 1,\\ \ldots,\ a_{n+1} = a_n + 1 \text{ ។}
\end{aligned}
\end{gathered}
\]
គេសង្កេតឃើញថា តួនីមួយៗនៃស៊ីត (ក្រៅពីតួទី $1$) ស្មើនឹងតួមុនបន្ទាប់បូកនឹងចំនួនថេរ (ស្មើនឹង $1$) ។ ចំនួន $1$ ហៅថាផលសងរួមនៃស៊ីត ។ ស៊ីត $1, 2, 3, 4, 5, \ldots$ ហៅថា \textbf{ស៊ីតនព្វន្ត}ដែលមានផលសងស្មើ $1$ ។
\end{example}

\begin{definition}
ស៊ីតនព្វន្ត គឺជាស៊ីតនៃចំនួនពិតដែលមានតួនីមួយៗ (ក្រៅពីតួទី $1$) ស្មើនឹងតួមុនបន្ទាប់បូកចំនួនថេរ $d$ មួយ ហៅថាផលសងរួម ។
\end{definition}

\begin{formulaBox}[ជាទូទៅ]
ផលសងរួមនៃស៊ីតនព្វន្តតាងដោយ $d$ កំណត់ដោយ
\[
d = a_2 - a_1 = a_3 - a_2 = \ldots = a_{n+1} - a_n \text{ ។}
\]
\end{formulaBox}

\section{តួទី \texorpdfstring{$n$}{n} នៃស៊ីតនព្វន្ត}

\begin{example}[ឧទាហរណ៍]
គេមានស៊ីត $1, 2, 3, 4, 5, \ldots$ គឺជាស៊ីតនព្វន្តដែលមានតួទី $1$ ៖ $a_1 = 1$ និងផលសងរួម $d = 1$ ។ គេបាន
\[
\begin{aligned}
a_1 &= 1,\ a_2 = 2 = 1 + 1 = a_1 + d,\ a_3 = 3 = 2 + 1\\
&= a_1 + 2d,
\end{aligned}
\]
\[
\begin{aligned}
a_4 &= 4 = 3 + 1 = a_1 + 3d,\ a_5 = 5 = 4 + 1\\
&= a_1 + 4d,\ \ldots
\end{aligned}
\]
គេសង្កេតឃើញ តួនីមួយៗនៃស៊ីត (ក្រៅពីតួទី $1$) ស្មើនឹងតួទី $1$ បូក $d, 2d, 3d, \ldots$ រៀងគ្នា ។
\end{example}

\begin{formulaBox}[ជាទូទៅ]
$(a_n)$ ជាស៊ីតនព្វន្តដែលមានតួទី $1$ $a_1$ និងផលសងរួម $d$ ។ តួទី $n$ នៃស៊ីតនព្វន្តកំណត់ដោយ
\[
a_n = a_1 + (n - 1) \times d \text{ ។}
\]
\end{formulaBox}

\begin{example}[លំហាត់គំរូ ១]
គេមានស៊ីតនព្វន្ត $2, 8, 14, 20, \ldots$ ។
\begin{enumerate}[a]
  \item គណនាតួទី $20$
  \item តើចំនួន $236$ ជាតួទីប៉ុន្មាននៃស៊ីត?
\end{enumerate}
\end{example}
\ans
តាងតួទីមួយ $a_1 = 2$ តួទីពីរ $a_2 = 8$ ។ គេបានផលសងរួម $d = a_2 - a_1 = 8 - 2 = 6$ ។
\begin{enumerate}[a]
  \item ដោយ $a_n = a_1 + (n - 1)d$ គេបានតួទី $20$ នៃស៊ីតនព្វន្ត
  \[ a_{20} = 2 + (20 - 1) \times 6 = 116 \]
  ដូច្នេះតួទី $20$ នៃស៊ីតគឺ $116$ ។
  \item តាង $a_n = 236$ គេបាន $a_n = 2 + (n - 1) \times 6$
  \[
\begin{gathered}
\begin{aligned}
236 &= 2 + (n - 1) \times 6 \iff 236 = 6n - 4\\ \iff 6n\\
&= 240 \Rightarrow n = 40
\end{aligned}
\end{gathered}
\]
  ដូច្នេះចំនួន $236$ ជាតួទី $40$ នៃស៊ីត ។
\end{enumerate}

\begin{example}[លំហាត់គំរូ ២]
តំណាងក្រុមហ៊ុនធានារ៉ាប់រងមួយទទួលប្រាក់បៀវត្សក្នុងឆ្នាំទី $1$ គឺ $30$ លានរៀល និងក្នុងឆ្នាំទី $7$ គឺ $60$ លានរៀល ។ ឧបមាថា តួលេខនៃប្រាក់បៀវត្សប្រចាំឆ្នាំរបស់គាត់បង្កើតបានជាស៊ីតនព្វន្ត ។ កំណត់ប្រាក់បៀវត្សរបស់គាត់ក្នុងឆ្នាំទី $10$ ។
\end{example}
\ans
គេមានប្រាក់បៀវត្សក្នុងឆ្នាំទី $1$ គឺ $a_1 = 30$ លានរៀល និងប្រាក់បៀវត្សក្នុងឆ្នាំទី $7$ $a_7 = 60$ លានរៀល នោះគឺ $a_7 = a_1 + (7 - 1)d$
\[
\begin{gathered}
60 = 30 + (7 - 1)d\\
\text{គេបាន } d = 5 \text{ លានរៀល}
\end{gathered}
\]
ដូច្នេះប្រាក់បៀវត្សក្នុងឆ្នាំទី $10$ គឺ
\[
\begin{gathered}
a_{10} = a_7 + (10 - 7)d,\\
a_{10} = 60 + 3 \times 5 = 75 \text{ លានរៀល ។}
\end{gathered}
\]

\begin{pro}
\begin{enumerate}[a]
  \item គេឱ្យស៊ីតនព្វន្តមួយដែល $a_1 = 2$ និង $a_2 = 9$ ។ កំណត់តួទី $5$ និងគណនា $a_6$ ។
  \item សៀវភៅមួយក្បាលមាន $100$ ទំព័រចុះលេខទំព័រពី $1$ ដល់ $100$ ។ រកចំនួនទំព័រដែលមានលេខខាងចុង $5$ ។
\end{enumerate}
\ans
\begin{enumerate}[a]
  \item $a_5 = 30$, $a_6 = 37$
  \item មាន $10$ ទំព័រ ។
\end{enumerate}
\end{pro}

\section{ផលបូក \texorpdfstring{$n$}{n} តួដំបូងនៃស៊ីតនព្វន្ត}

គេមានស៊ីតនព្វន្ត $a_1, a_2, a_3, a_4, a_5, a_6$ ។ $a_1$ និង $a_6$ ហៅថា \textbf{តួចុង} ។ $a_2$ និង $a_5$ ហើយ $a_3$ និង $a_4$ ហៅថា \textbf{តួនៅស្មើចម្ងាយពីតួចុង} ។

តាមរូបមន្តតួទី $n$ នៃស៊ីតនព្វន្ត ៖ $a_n = a_1 + (n - 1)d$ គេបាន
\[
\begin{aligned}
a_1 + a_6 &= a_1 + a_1 + 5d = 2a_1 + 5d &\quad (1)\\
a_2 + a_5 &= a_1 + d + a_1 + 4d = 2a_1 + 5d &\quad (2)\\
a_3 + a_4 &= a_1 + 2d + a_1 + 3d = 2a_1 + 5d &\quad (3)
\end{aligned}
\]
គេបាន $a_1 + a_6 = a_2 + a_5 = a_3 + a_4$ ។

ជាទូទៅ បើគេមានស៊ីតនព្វន្ត $a_1, a_2, a_3, \ldots, a_{n-2}, a_{n-1}, a_n$ គេបានផលបូកតួចុងស្មើនឹងផលបូកតួស្មើចម្ងាយពីតួចុង
\[
a_1 + a_n = a_2 + a_{n-1} = a_3 + a_{n-2} = \ldots
\]
គេប្រើលក្ខណៈនេះដើម្បីគណនាផលបូកតួនៃស៊ីតនព្វន្ត
\[
\begin{aligned}
S_n &= a_1 + a_2 + a_3 + \ldots + a_{n-1} + a_n\\
S_n &= a_n + a_{n-1} + \ldots + a_2 + a_1
\end{aligned}
\]
បូកអង្គនិងអង្គ គេបាន
\[
\begin{aligned}
2S_n &= (a_1 + a_n) + (a_2 + a_{n-1}) + \ldots\\
&\quad + (a_{n-1} + a_2) + (a_1 + a_n)\\
&= \underbrace{(a_1 + a_n) + (a_1 + a_n) + \ldots + (a_1 + a_n)}_{n \text{ តួ}}
\end{aligned}
\]
\[
\begin{gathered}
2S_n = (a_1 + a_n)n \quad \text{នោះ}\\
S_n = \frac{n}{2}(a_1 + a_n) \text{ ។}
\end{gathered}
\]

\begin{formulaBox}[ជាទូទៅ]
ផលបូក $n$ តួដំបូងនៃស៊ីតនព្វន្ត ដែលមាន $a_1$ ជាតួទី $1$ និង $a_n$ ជាតួទី $n$ កំណត់ដោយ
\[
\begin{gathered}
S_n = \frac{n}{2}(a_1 + a_n) \quad \text{ឬ}\\
S_n = \frac{n\left(2a_1 + (n - 1)d\right)}{2} \text{ ។}
\end{gathered}
\]
\end{formulaBox}

\begin{example}[លំហាត់គំរូ ១]
គេមានស៊ីតនព្វន្ត $51, 47, 43, \ldots$ ។ កំណត់តម្លៃ $n$ ដែលគេឱ្យផលបូក $n$ តួដំបូង $S_n$ មានតម្លៃអតិបរមា និងកំណត់តម្លៃនៃ $S_n$ ។
\end{example}
\ans
តាង $a_n$ ជាតួទូទៅនៃស៊ីតនព្វន្ត ។ ដោយផលសងរួម $d = -4$ និង $a_1 = 51$ ។ គេបាន
\[ a_n = a_1 + (n - 1)d = 51 + (n - 1)(-4) = 55 - 4n \]
នាំឱ្យ $a_n > 0$ ចំពោះ $n = 1, 2, \ldots, 13$ និង $a_n < 0$ ចំពោះ $n = 14, 15, \ldots$ ។

ដូច្នេះ $S_n$ អតិបរមានៅពេលដែល $n = 13$ និង
\[ S_{13} = \frac{51 + (55 - 4 \times 13)}{2} \times 13 = 351 \text{ ។} \]

\begin{example}[លំហាត់គំរូ ២]
គេរៀបឥដ្ឋទីធ្លាមុខផ្ទះមួយដែលមានរាងជាចតុកោណព្នាយ (មើលរូប) ។ ទីធ្លានោះមានឥដ្ឋ $18$ ជួរ ។ ជួរទី $1$ មានឥដ្ឋ $14$ ដុំ ហើយជួរបន្តបន្ទាប់មានចំនួនឥដ្ឋលើសជួរមុន $1$ ដុំ រហូតដល់ជួរទី $18$ មាន $31$ ដុំ ។ តើគេត្រូវចំណាយឥដ្ឋប៉ុន្មានដុំដើម្បីរៀបឱ្យពេញទីធ្លានោះ?

\begin{center}
\begin{tikzpicture}[scale=0.6]
  \fill[orange!80!red] (0,0) -- (8,0) -- (6.5,-1.6) -- (1.5,-1.6) -- cycle;
  \foreach \y in {-0.2,-0.4,...,-1.4} \draw[white,thin] (0.19*-\y*8/1.6+0,\y) -- (8-0.19*-\y*8/1.6,\y);
  \node[above] at (4,0) {$31$ ដុំ};
  \node[below] at (4,-1.6) {$14$ ដុំ};
\end{tikzpicture}
\end{center}
\end{example}
\ans
ចំនួនឥដ្ឋតាមជួរដែលរៀបទីធ្លាមុខផ្ទះនេះបង្កើតបានជាស៊ីតនព្វន្តដែលមានតួទី $1$ គឺ $a_1 = 14$, $a_{18} = 31$ និងមាន $18$ ជួរ ។ តាមរូបមន្តផលបូកតួនៃស៊ីតនព្វន្ត គេបាន
\[ S_{18} = \frac{(a_1 + a_{18}) \times 18}{2} = \frac{(14 + 31) \times 18}{2} = 405 \]
ការរៀបចំទីធ្លាមុខផ្ទះ គេត្រូវចំណាយឥដ្ឋអស់ $405$ ដុំ ។

\begin{example}[លំហាត់គំរូ ៣]
គេមានការងាររដូវក្តៅ $2$ សម្រាប់រយៈពេល $3$ ខែឬ $12$ សប្តាហ៍ ។ ការងារ $A$ ទទួលបានប្រាក់បៀវត្ស $400\,000$ រៀលក្នុងមួយខែ ហើយដំឡើង $100\,000$ រៀលរៀងរាល់ខែ ។ ការងារ $B$ ទទួលប្រាក់បៀវត្ស $100\,000$ រៀលក្នុងមួយសប្តាហ៍ហើយដំឡើង $5\,000$ រៀលរៀងរាល់សប្តាហ៍ ។ តើគេគួរជ្រើសរើសយកការងារណាមួយប្រសើរជាង?
\end{example}
\ans
ចំពោះការងារ $A$ ប្រាក់បៀវត្សសរុបសម្រាប់រយៈពេល $3$ ខែគឺ
\[ S_3 = 400\,000 + 500\,000 + 600\,000 = 1\,500\,000 \text{ រៀល} \]
ការងារ $B$ ប្រាក់បៀវត្សគឺ $100\,000, 105\,000, 110\,000, \ldots, a_{12}$ ។ នេះជាស៊ីតនព្វន្តដែលមានតួទី $1$ គឺ $a_1 = 100\,000$ និងផលសងរួម $d = 5\,000$ ។ តាមរូបមន្ត ប្រាក់បៀវត្សនៅសប្តាហ៍ទី $12$ គឺ
\[
\begin{aligned}
a_{12} &= a_1 + (12 - 1)d = 100\,000 + 11 \times 5\,000\\
&= 155\,000 \text{ ។}
\end{aligned}
\]
ចំពោះការងារ $B$ ប្រាក់បៀវត្សសរុបរយៈពេល $3$ ខែ ឬ $12$ សប្តាហ៍គឺ
\[
\begin{aligned}
S_B &= 100\,000 + 105\,000 + 110\,000 + \ldots + 155\,000\\
&= \frac{(100\,000 + 155\,000) \times 12}{2}\\
&= 255\,000 \times 6 = 1\,530\,000 \text{ រៀល ។}
\end{aligned}
\]
ប្រាក់បានពីការងារ $B$ ច្រើនជាងបានពីការងារ $A$ ដូច្នេះគេរើសយកការងារ $B$ ។

\begin{pro}
\begin{enumerate}[1]
  \item \begin{enumerate}[a]
    \item គណនាផលបូក $(-6) + (-1) + 4 + 9 + \ldots + 64$ ។
    \item គណនាផលបូក $25$ តួដំបូងនៃស៊ីតនព្វន្ត $2, 7, 12, 17, \ldots$
  \end{enumerate}
  \item គេមានការងាររដូវក្តៅ $2$ សម្រាប់រយៈពេល $3$ ខែឬ $12$ សប្តាហ៍ ។ ការងារ $A$ ទទួលបានប្រាក់បៀវត្ស $325\,000$ រៀលក្នុងមួយខែ ហើយដំឡើង $100\,000$ រៀលរៀងរាល់ខែ ។ ការងារ $B$ ទទួលបានប្រាក់បៀវត្ស $80\,000$ រៀលក្នុងមួយសប្តាហ៍ ហើយដំឡើង $3\,000$ រៀលរៀងរាល់សប្តាហ៍ ។ តើគេគួរជ្រើសរើសយកការងារណាមួយប្រសើរជាង?
\end{enumerate}
\ans
\begin{enumerate}[1]
  \item ក. $S_{15} = 435$ \quad ខ. $S_{25} = 1550$ ។
  \item គេរើសយកការងារ $A$ ។
\end{enumerate}
\end{pro}

\section*{លំហាត់អនុវត្តន៍}
\begin{enumerate}[1]
  \item សរសេររូបមន្តតួទី $n$ នៃស៊ីតនព្វន្តខាងក្រោម ៖
  \begin{tasks}(2)
    \task $4, 15, 26, 37, \ldots$
    \task $5, 1, -3, -7, -11, \ldots$
    \task $-2, -9, -16, -23, \ldots$
    \task $7, 12, 17, 22, 27, \ldots$ ។
  \end{tasks}
  \item សរសេរប្រាំតួដំបូងនៃស៊ីតនព្វន្តដែលកំណត់ដោយតួ
  \begin{tasks}(2)
    \task $a_1 = 3$, $a_2 = 10$
    \task $a_n = -7 + 3n$
  \end{tasks}
  \item គេមានស៊ីតនព្វន្ត $-8, -3, 2, 7, \ldots$ ។
  \begin{tasks}(2)
    \task គណនា $a_{17}$ និង $a_{43}$
    \task តើតួទីប៉ុន្មាននៃស៊ីតស្មើនឹង $322$ ។
  \end{tasks}
  \item កំណត់ចំនួនតួនៃស៊ីតនព្វន្តខាងក្រោម ៖
  \begin{tasks}(2)
    \task $18, 13, 8, \ldots, -102$
    \task $-9, -4, 1, \ldots, 171$ ។
  \end{tasks}
  \item ក្នុងចំណោមស៊ីតខាងក្រោម តើស៊ីតណាជាស៊ីតនព្វន្ត?
  \begin{tasks}(2)
    \task $a_n = 13 + 6(n - 1)$
    \task $a_1 = 3$, $a_{n+1} = a_n - \dfrac{1}{4}$ ។
  \end{tasks}
  \item គេឱ្យ $(a_n)$ ជាស៊ីតនព្វន្ត បើគេដឹងថា
  \begin{enumerate}[a]
    \item $a_{35} = 9$, $a_{45} = 17$ ។ គណនា $a_{40}$ ។
    \item $a_7 = \dfrac{7}{2}$, $a_{13} = \dfrac{13}{2}$ ។ គណនា $a_1$ ។
    \item $a_2 = -12$, $S_{12} = 18$ ។ គណនា $a_1$, $d$ និង $U_6$ ។
  \end{enumerate}
  \item គេឱ្យស៊ីតនព្វន្ត $2, 7, 12, \ldots$ និងស៊ីតនព្វន្ត $2, 5, 8, \ldots$ ។ តើក្នុងចំណោម $121$ តួនៃស៊ីតទាំងពីរនេះមានប៉ុន្មានតួដែលមានតម្លៃស្មើៗគ្នា?
  \item បង្ហាញថាស៊ីត $(a_n)$ ដែលមានតួទី $n$ កំណត់ដោយ $a_n = \dfrac{2n+3}{4}$ ជាស៊ីតនព្វន្តកើន ។
  \item គេឱ្យប្រាំចំនួនជាស៊ីតនព្វន្ត ។ គណនាចំនួនទាំងនោះបើគេដឹងថាផលបូកវាស្មើនឹង $125$ ហើយផលបូកពីរតួដំបូងស្មើនឹង $38$ ។
  \item បើផលបូក $9$ តួដំបូងនៃស៊ីតស្មើនឹង $162$ ហើយផលបូក $12$ តួដំបូងនៃស៊ីតនព្វន្តស្មើនឹង $288$ ។ គណនាផលបូក $30$ តួដំបូងនៃស៊ីតនព្វន្ត ។
  \item បីចំនួនតគ្នាជាស៊ីតនព្វន្តដោយដឹងថាផលបូកនៃបីចំនួននេះស្មើនឹង $36$ និងផលគុណវាស្មើនឹង $1428$ ។
  \item គណនាផលបូកស៊ីតនព្វន្តខាងក្រោម ៖
  \begin{tasks}(2)
    \task $2 + 4 + 6 + \ldots + 146$
    \task $100 + 95 + 90 + 85 + \ldots - 20$
    \task $-193 - 189 - 185 - \ldots - 21 - 17$
    \task $5\frac{1}{4} + 4\frac{1}{2} + 3\frac{3}{4} + \ldots - 3$ ។
  \end{tasks}
  \item អ្នកជិះកង់ម្នាក់ធ្វើដំណើរតាមផ្លូវ $A$ បានចម្ងាយ $15$ គីឡូម៉ែត្រក្នុង $1$ ថ្ងៃ ។ ប្រាំបីថ្ងៃក្រោយមក ទើបអ្នកជិះកង់ទីពីរចាប់ផ្តើមធ្វើដំណើរតាមផ្លូវ $A$ ដូចគ្នាដោយថ្ងៃទីមួយ គាត់ជិះកង់បានចម្ងាយ $14$ គីឡូម៉ែត្រក្នុង $1$ ថ្ងៃ ហើយថ្ងៃបន្តបន្ទាប់មកទៀតគាត់ត្រូវជិះកើនបានចម្ងាយ $2$ គីឡូម៉ែត្រក្នុង $1$ ថ្ងៃជារៀងរាល់ថ្ងៃដើម្បីឱ្យទាន់អ្នកជិះកង់ទីមួយ $A$ ។ តើអ្នកជិះកង់ទីពីរត្រូវចំណាយរយៈពេលអស់ប៉ុន្មានថ្ងៃ ដើម្បីជិះទៅកាន់អ្នកជិះកង់ទីមួយ?
  \item សាលាប្រជុំមួយមានរាងចតុកោណព្នាយ ហើយតាមជ្រុងទ្រេតនៃសាលាប្រជុំនោះ គេរៀបកៅអីជា $15$ ជួរដែលជួរមុខគេបង្អស់រៀបបាន $25$ កៅអី ហើយជួរទី $2$ រៀបបាន $27$ កៅអី ជួរទី $3$ រៀបបាន $29$ កៅអី ។ គេរៀបកៅអីតាមលំនាំខាងលើរហូតដល់ជួរខាងក្រោយ ។ គណនាកៅអីសម្រាប់អង្គុយទាំងអស់ដែលមានក្នុងសាលប្រជុំ ។
\end{enumerate}
